% front.tex (국문 v13): 제목 · 저자 · 초록 · 핵심어(블록 1), 자료 가용성(블록 3), 뒷부분(블록 4).
\providecommand{\frontblock}{0}
\ifnum\frontblock=1
\title[물리 잔차 기반 기계학습의 연도별 활동층 두께 예측]{물리 잔차 기반 기계학습 워크플로를 활용한 영구동토 활동층 두께 예측과 지역별 편차 분석}

\author[1]{\fnm{백승원}}
\author[2]{\fnm{김승희}}
\author[2]{\fnm{정윤택}}
\author*[1]{\fnm{조용채}}\email{[교신 저자 전자우편]}
\affil*[1]{\orgdiv{에너지시스템공학부}, \orgname{서울대학교}, \orgaddress{\city{서울}, \country{대한민국}}}
\affil[2]{\orgdiv{원격탐사빙권정보센터}, \orgname{극지연구소}, \orgaddress{\city{인천}, \country{대한민국}}}

\abstract{활동층 두께(ALT)는 해마다 달라지지만, 기존 격자 ALT 제품은 대부분 여러 해의 조건으로 학습되었고 같은 위치 · 같은 해의 관측과 비교된 적이 드물다. 이 연구는 알래스카, 레나 델타, 캐나다 북서부의 지점-연도 ALT 관측 \nncyAll{}개(2010--2024년)에 관측일까지 누적한 ERA5-Land 융해 도일을 결합하고, 지역 라벨로 재보정한 Stefan 모델의 잔차를 gradient boosting 으로 학습하는 물리 잔차 워크플로를 검증하였다. 공간 블록 교차검증에서 물리 잔차 모델은 Stefan 모델의 오차를 레나 델타에서 \ndrefLeblkRoabs{} cm, 캐나다에서 \ndrefCablkRoabs{} cm 줄였고 알래스카에서는 같은 수준이었다. 이득은 토양 · 지형 공변량을 요구하였고 지역 Stefan 모델의 오차가 큰 곳에서 가장 컸다. 레나 델타를 학습에서 제외하면 직접 기계학습의 오차는 Stefan 모델보다 \ndrefLeregDzabs{} cm 컸으나 물리 잔차 모델은 Stefan 모델 수준을 유지하였고, 지역 라벨 25개로 그 오차는 \ntlrmseLeRszero{} cm에서 \ntlrmseLeRstwfive{} cm로 줄었다. 같은 셀 · 같은 해 비교에서 레나 델타의 오차(\nccLeRoall{} cm)는 얇은 활동층을 나타내지 못하는 공개 제품 5종(\nccLeCcifiveall--\nccLeWeiall{} cm)보다 낮았다. 이 워크플로는 세 지역의 2011--2024년 연도별 1 km ALT 지도를 적용 가능 영역 · 예측 구간과 함께 제공한다.}

\keywords{영구동토, 활동층 두께, Stefan 모델, 물리 기반 기계학습, 공간 교차검증, ERA5-Land}
\fi

\ifnum\frontblock=3
\bmhead{자료 가용성}
모든 입력 자료는 공개 자료이다. ALT 관측: ABoVE 토양 수분 · 활동층 두께 v2(ORNL DAAC, \url{https://doi.org/10.3334/ORNLDAAC/2369}), ALLena 융해 깊이 자료(PANGAEA, \url{https://doi.org/10.1594/PANGAEA.973813}), GTN-P CALM(PANGAEA, \url{https://doi.org/10.1594/PANGAEA.972777}). 기후: ERA5-Land 월평균(Copernicus Climate Data Store). 공변량: Copernicus DEM GLO-30, SoilGrids 2.0, ESA CCI 영구동토 확률, MODIS MOD13Q1 · MOD10A2 v6.1(AppEEARS). 제품: Wei 등(\url{https://doi.org/10.5281/zenodo.21667583}), Liu 등(\url{https://doi.org/10.5281/zenodo.10070610}), ESA CCI ALT v4 · v5(CEDA), Aalto 등(\url{https://zenodo.org/records/5003804}), Yi와 Kimball(\url{https://doi.org/10.3334/ORNLDAAC/1760}). 지점-연도 라벨 표, 연도별 기후 입력 표, 정적 공변량 표, 모든 방법의 겹 밖 예측, 알래스카 · 레나 델타 · 캐나다 북서부의 연도별 1 km 지도(2011--2024년)는 게재 시 공개 저장소에 둔다[저장소 · DOI].
\fi

\ifnum\frontblock=4
\bmhead{감사의 글}
[연구비 · 감사의 글.]

\bmhead{저자 기여}
백승원은 연구를 설계하고 자료를 구축하였으며 분석을 수행하고 원고를 작성하였다. 김승희와 정윤택은 원격탐사 공변량과 영구동토 현장 자료의 해석에 기여하고 원고를 검토하였다. 조용채는 연구를 지도하고 원고를 검토하였다. 모든 저자가 최종 원고를 승인하였다.

\bmhead{추가 정보}
\textbf{보충 자료} 보충 방법 S1--S5, 보충 표 S1--S7, 보충 그림 S1--S10을 함께 제공한다.

\textbf{이해 상충} 저자들은 이해 상충이 없음을 밝힌다.
\fi
